% GENERATED FILE. Edit canonical lessons, then run npm run generate:books.
% canonical-source-hash: fnv1a64:7c3e577c835ed02f
% canonical-lessons: TE-C183-vantagadi, TE-C183-padakagadi, TE-C183-balci, TE-C183-muta, TE-C183-banali, TE-R183-first-pass-words-from-chapters-175-178, TE-R183-second-pass-words-from-chapters-179-183

\chapter{Kitchen, Bedroom, Bucket, Lid, Frying pan}
\label{ch:te-kitchen-bedroom-bucket-lid-frying-pan}
\hlchaptermodality{183}

\begin{chapteropening}
\textbf{By the end of this chapter:} \emph{I can say a kitchen, a bedroom, a bucket, a lid and a frying pan.}

Complete the last lesson of chapter 183: I can say a kitchen, a bedroom, a bucket, a lid and a frying pan.
\end{chapteropening}

\section[\texorpdfstring{\te{వంటగది} (vaṇṭagadi)}{వంటగది (vaṇṭagadi)}]{\te{వంటగది} (vaṇṭagadi) — a kitchen}

\label{lesson:TE-C183-vantagadi}



\begin{quote}
Before the new one: say the Telugu for plain, then the Telugu for extra.
\end{quote}

\subsection*{You'll want to know: \te{వంటగది}}
\textbf{\te{వంటగది}} — \emph{vaṇṭagadi} — \textquotedblleft{}a kitchen\textquotedblright{}.

\begin{grammarlens}[title={What you've built}]
One more word for this chapter.
\end{grammarlens}

\subsection*{Your turn}
\begin{itemize}
\raggedright
  \item \emph{Say it:} \emph{vaṇṭagadi}
  \item \emph{Say it:} \emph{vaṇṭagadi}, once more
  \item \emph{Say it:} say \emph{vaṇṭagadi}
  \item \emph{Recall:} say the Telugu for plain, then the Telugu for extra, then say \emph{vaṇṭagadi} again
\end{itemize}

\begin{tcolorbox}[breakable,colback=teal!4,colframe=teal!35!black,title={Before you move on}]
What does \te{వంటగది} mean? (\textquotedblleft{}A kitchen\textquotedblright{}.) Say it once more.
\end{tcolorbox}
\section[\texorpdfstring{\te{పడకగది} (paḍakagadi)}{పడకగది (paḍakagadi)}]{\te{పడకగది} (paḍakagadi) — a bedroom}

\label{lesson:TE-C183-padakagadi}



\begin{quote}
Before the new one: say the Telugu for extra, then the Telugu for a kitchen.
\end{quote}

\subsection*{You'll want to know: \te{పడకగది}}
\textbf{\te{పడకగది}} — \emph{paḍakagadi} — \textquotedblleft{}a bedroom\textquotedblright{}.

\begin{grammarlens}[title={What you've built}]
One more word for this chapter.
\end{grammarlens}

\subsection*{Your turn}
\begin{itemize}
\raggedright
  \item \emph{Say it:} \emph{paḍakagadi}
  \item \emph{Say it:} \emph{paḍakagadi}, once more
  \item \emph{Say it:} say \emph{paḍakagadi}
  \item \emph{Recall:} say the Telugu for extra, then the Telugu for a kitchen, then say \emph{paḍakagadi} again
\end{itemize}

\begin{tcolorbox}[breakable,colback=teal!4,colframe=teal!35!black,title={Before you move on}]
What does \te{పడకగది} mean? (\textquotedblleft{}A bedroom\textquotedblright{}.) Say it once more.
\end{tcolorbox}
\section[\texorpdfstring{\te{బాల్చీ} (bālcī)}{బాల్చీ (bālcī)}]{\te{బాల్చీ} (bālcī) — a bucket}

\label{lesson:TE-C183-balci}



\begin{quote}
Before the new one: say the Telugu for a kitchen, then the Telugu for a bedroom.
\end{quote}

\subsection*{You'll want to know: \te{బాల్చీ}}
\textbf{\te{బాల్చీ}} — \emph{bālcī} — \textquotedblleft{}a bucket\textquotedblright{}.

\begin{grammarlens}[title={What you've built}]
One more word for this chapter.
\end{grammarlens}

\subsection*{Your turn}
\begin{itemize}
\raggedright
  \item \emph{Say it:} \emph{bālcī}
  \item \emph{Say it:} \emph{bālcī}, once more
  \item \emph{Say it:} say \emph{bālcī}
  \item \emph{Recall:} say the Telugu for a kitchen, then the Telugu for a bedroom, then say \emph{bālcī} again
\end{itemize}

\begin{tcolorbox}[breakable,colback=teal!4,colframe=teal!35!black,title={Before you move on}]
What does \te{బాల్చీ} mean? (\textquotedblleft{}A bucket\textquotedblright{}.) Say it once more.
\end{tcolorbox}
\section[\texorpdfstring{\te{మూత} (mūta)}{మూత (mūta)}]{\te{మూత} (mūta) — a lid}

\label{lesson:TE-C183-muta}



\begin{quote}
Before the new one: say the Telugu for a bedroom, then the Telugu for a bucket.
\end{quote}

\subsection*{You'll want to know: \te{మూత}}
\textbf{\te{మూత}} — \emph{mūta} — \textquotedblleft{}a lid\textquotedblright{}.

\begin{grammarlens}[title={What you've built}]
One more word for this chapter.
\end{grammarlens}

\subsection*{Your turn}
\begin{itemize}
\raggedright
  \item \emph{Say it:} \emph{mūta}
  \item \emph{Say it:} \emph{mūta}, once more
  \item \emph{Say it:} say \emph{mūta}
  \item \emph{Recall:} say the Telugu for a bedroom, then the Telugu for a bucket, then say \emph{mūta} again
\end{itemize}

\begin{tcolorbox}[breakable,colback=teal!4,colframe=teal!35!black,title={Before you move on}]
What does \te{మూత} mean? (\textquotedblleft{}A lid\textquotedblright{}.) Say it once more.
\end{tcolorbox}
\section[\texorpdfstring{\te{బాణలి} (bāṇali)}{బాణలి (bāṇali)}]{\te{బాణలి} (bāṇali) — a frying pan, a wok}

\label{lesson:TE-C183-banali}



\begin{quote}
Before the new one: say the Telugu for a bucket, then the Telugu for a lid.
\end{quote}

\subsection*{You'll want to know: \te{బాణలి}}
\textbf{\te{బాణలి}} — \emph{bāṇali} — \textquotedblleft{}a frying pan, a wok\textquotedblright{}.

\begin{grammarlens}[title={What you've built}]
One more word for this chapter.
\end{grammarlens}

\subsection*{Your turn}
\begin{itemize}
\raggedright
  \item \emph{Say it:} \emph{bāṇali}
  \item \emph{Say it:} \emph{bāṇali}, once more
  \item \emph{Say it:} say \emph{bāṇali}
  \item \emph{Recall:} say the Telugu for a bucket, then the Telugu for a lid, then say \emph{bāṇali} again
\end{itemize}

\begin{tcolorbox}[breakable,colback=teal!4,colframe=teal!35!black,title={Before you move on}]
What does \te{బాణలి} mean? (\textquotedblleft{}A frying pan, a wok\textquotedblright{}.) Say it once more.
\end{tcolorbox}
\section[(dialogue)]{Words from chapters 175-178 — a first pass}

\label{lesson:TE-R183-first-pass-words-from-chapters-175-178}



\begin{quote}
Say the words for a lid and a frying pan.
\end{quote}

\subsection*{Your turn}
Say these aloud:

\begin{itemize}
\raggedright
  \item a crowd, the world, an earthquake, dust, smoke
  \item otherwise, even so, through, about, without
  \item if by chance, instead of, according to, actually, anyway
  \item on purpose, by mistake, never mind, a condition, an excuse
\end{itemize}

\begin{tcolorbox}[breakable,colback=teal!4,colframe=teal!35!black,title={Before you move on}]
A crowd? (\textbf{\te{గుంపు}}, \emph{gumpu}.) The world? (\textbf{\te{ప్రపంచం}}, \emph{prapañcaṁ}.) An earthquake? (\textbf{\te{భూకంపం}}, \emph{bhūkampaṁ}.) Dust? (\textbf{\te{దుమ్ము}}, \emph{dummu}.) Smoke? (\textbf{\te{పొగ}}, \emph{poga}.) Otherwise? (\textbf{\te{లేకపోతే}}, \emph{lēkapōtē}.) Even so? (\textbf{\te{అయినా}}, \emph{ayinā}.) Through? (\textbf{\te{ద్వారా}}, \emph{dvārā}.) About? (\textbf{\te{గురించి}}, \emph{guriñci}.) Without? (\textbf{\te{లేకుండా}}, \emph{lēkuṇḍā}.) If by chance? (\textbf{\te{ఒకవేళ}}, \emph{okavēḷa}.) Instead of? (\textbf{\te{బదులు}}, \emph{badulu}.) According to? (\textbf{\te{ప్రకారం}}, \emph{prakāraṁ}.) Actually? (\textbf{\te{అసలు}}, \emph{asalu}.) Anyway? (\textbf{\te{ఏదేమైనా}}, \emph{ēdēmainā}.) On purpose? (\textbf{\te{కావాలని}}, \emph{kāvālani}.) By mistake? (\textbf{\te{పొరపాటున}}, \emph{porapāṭuna}.) Never mind? (\textbf{\te{పోనీ}}, \emph{pōnī}.) A condition? (\textbf{\te{షరతు}}, \emph{ṣaratu}.) An excuse? (\textbf{\te{సాకు}}, \emph{sāku}.)
\end{tcolorbox}
\section[(dialogue)]{Words from chapters 179-183 — a second pass}

\label{lesson:TE-R183-second-pass-words-from-chapters-179-183}



\begin{quote}
Say the words for delicious and bland.
\end{quote}

\subsection*{Your turn}
Say these aloud:

\begin{itemize}
\raggedright
  \item delicious, bland, fresh, expensive, cheap
  \item main, total, every, pink, orange
  \item great, a smell, interesting, wonderful, ordinary
  \item silence, durable, fake, plain, extra
  \item a kitchen, a bedroom, a bucket, a lid, a frying pan
\end{itemize}

\begin{tcolorbox}[breakable,colback=teal!4,colframe=teal!35!black,title={Before you move on}]
Delicious? (\textbf{\te{కమ్మని}}, \emph{kammani}.) Bland? (\textbf{\te{చప్పని}}, \emph{cappani}.) Fresh? (\textbf{\te{తాజా}}, \emph{tājā}.) Expensive? (\textbf{\te{ఖరీదైన}}, \emph{kharīdaina}.) Cheap? (\textbf{\te{చౌక}}, \emph{cauka}.) Main? (\textbf{\te{ప్రధాన}}, \emph{pradhāna}.) Total? (\textbf{\te{మొత్తం}}, \emph{mottaṁ}.) Every? (\textbf{\te{ప్రతి}}, \emph{prati}.) Pink? (\textbf{\te{గులాబీ} \te{రంగు}}, \emph{gulābī raṅgu}.) Orange? (\textbf{\te{నారింజ} \te{రంగు}}, \emph{nāriñja raṅgu}.) Great? (\textbf{\te{గొప్ప}}, \emph{goppa}.) A smell? (\textbf{\te{వాసన}}, \emph{vāsana}.) Interesting? (\textbf{\te{ఆసక్తికరమైన}}, \emph{āsaktikaramaina}.) Wonderful? (\textbf{\te{అద్భుతమైన}}, \emph{adbhutamaina}.) Ordinary? (\textbf{\te{మామూలు}}, \emph{māmūlu}.) Silence? (\textbf{\te{నిశ్శబ్దం}}, \emph{niśśabdaṁ}.) Durable? (\textbf{\te{మన్నికైన}}, \emph{mannikaina}.) Fake? (\textbf{\te{నకిలీ}}, \emph{nakilī}.) Plain? (\textbf{\te{సాదా}}, \emph{sādā}.) Extra? (\textbf{\te{అదనపు}}, \emph{adanapu}.) A kitchen? (\textbf{\te{వంటగది}}, \emph{vaṇṭagadi}.) A bedroom? (\textbf{\te{పడకగది}}, \emph{paḍakagadi}.) A bucket? (\textbf{\te{బాల్చీ}}, \emph{bālcī}.) A lid? (\textbf{\te{మూత}}, \emph{mūta}.) A frying pan? (\textbf{\te{బాణలి}}, \emph{bāṇali}.)
\end{tcolorbox}
